\documentclass[11pt]{amsart}
\usepackage[T1]{fontenc}
\usepackage[utf8]{inputenc}
\usepackage{lmodern}
\usepackage[margin=1in]{geometry}
\usepackage{amsmath,amssymb,amsthm,mathtools}
\usepackage[hidelinks]{hyperref}
\allowdisplaybreaks

\newtheorem{theorem}{Theorem}
\newtheorem{lemma}[theorem]{Lemma}
\newtheorem{proposition}[theorem]{Proposition}
\theoremstyle{remark}
\newtheorem{remark}[theorem]{Remark}

\newcommand{\RH}{\mathrm{RH}}
\newcommand{\Ree}{\operatorname{Re}}
\newcommand{\ord}{\operatorname{ord}}

\title{M\"obius--Centroid Closure of the Left Half of the Critical Strip}
\author{}
\date{}

\begin{document}
\maketitle

\section{Ordinal zero sequence and multiplicity blocks}

Let
\[
\rho_k=\beta_k+i\gamma_k,\qquad k=1,2,3,\ldots,
\]
enumerate all nontrivial zeros of the Riemann zeta function in the upper half-plane, counted with multiplicity and ordered by nondecreasing positive ordinate.  Thus the running index is always
\[
k=1,2,3,\ldots,N,
\]
and $N$ denotes only the final ordinal index of a finite truncation.

Equal ordinates are kept as repeated ordinal positions.  If a positive ordinate $\Gamma_j$ occupies the block
\[
I_j=\{k_j,k_j+1,\ldots,k_j+m_j-1\},
\]
then
\[
\gamma_k=\Gamma_j\qquad(k\in I_j).
\]
The length $m_j$ contains every upper-half-plane zero having that ordinate, with multiplicity.  In particular, multiplicity on the critical line and, off the critical line, the functional-equation partner with the same positive ordinate both occupy genuine ordinal positions and shift every later index.

\section{The holomorphic numerator family}

For $u\in\mathbb C$ with $\Im u>0$, define
\begin{equation}\label{eq:Zfamily}
Z_\gamma(s,u)
:=
\sum_{k\ge1}\frac{e^{iu\gamma_k}}{k^s}.
\end{equation}
Since
\[
|e^{iu\gamma_k}|=e^{-(\Im u)\gamma_k},
\]
and $\gamma_k\asymp k/\log k$ with multiplicity counted in the ordinal enumeration, the exponential factor dominates every power of $k$.  Hence \eqref{eq:Zfamily} converges absolutely and locally uniformly in $s$ on every compact subset of $\mathbb C$.

\begin{proposition}\label{prop:entire}
For every $u$ with $\Im u>0$, the function $s\mapsto Z_\gamma(s,u)$ is entire.
\end{proposition}

Grouping by equal-ordinate blocks gives
\begin{equation}\label{eq:blockExpansion}
Z_\gamma(s,u)
=
\sum_{j\ge1} A_j(s)e^{iu\Gamma_j},
\qquad
A_j(s):=\sum_{k\in I_j}k^{-s}.
\end{equation}

\begin{lemma}[Non-cancellation for large imaginary frequency]\label{lem:noncancel}
For every fixed $s_0\in\mathbb C$, there exists $Y(s_0)>0$ such that
\[
Z_\gamma(s_0,iy)\neq0
\qquad (y>Y(s_0)).
\]
\end{lemma}

\begin{proof}
For a block $I_j=\{k_j,\ldots,k_j+m_j-1\}$,
\[
A_j(s_0)
=
k_j^{-s_0}
\sum_{\ell=0}^{m_j-1}
\left(1+\frac{\ell}{k_j}\right)^{-s_0}.
\]
The standard zero-counting bounds imply $m_j=O(\log \Gamma_j)$, while $k_j\asymp \Gamma_j\log\Gamma_j$, and hence
\[
\frac{m_j}{k_j}\longrightarrow0.
\]
Therefore, for fixed $s_0$,
\[
A_j(s_0)
=
m_jk_j^{-s_0}
\left(1+O_{s_0}\!\left(\frac{m_j}{k_j}\right)\right),
\]
so $A_j(s_0)\neq0$ for every sufficiently large $j$.  Let $j_0$ be the first block for which $A_{j_0}(s_0)\neq0$.  From \eqref{eq:blockExpansion},
\[
e^{y\Gamma_{j_0}}Z_\gamma(s_0,iy)
=
A_{j_0}(s_0)
+
\sum_{j>j_0}A_j(s_0)e^{-y(\Gamma_j-\Gamma_{j_0})}.
\]
Absolute convergence for every $y>0$ and strict increase of the distinct ordinates imply that the tail tends to $0$ as $y\to\infty$.  Hence
\[
e^{y\Gamma_{j_0}}Z_\gamma(s_0,iy)
\longrightarrow
A_{j_0}(s_0)\neq0,
\]
which proves the claim.
\end{proof}

\section{M\"obius transform and the three-factor identity}

Define
\[
a_k(u):=e^{iu\gamma_k},
\]
and the M\"obius-transformed coefficients
\begin{equation}\label{eq:bn}
b_n(u)
:=
\sum_{d\mid n}\frac{\mu(d)}{d}\,a_{n/d}(u).
\end{equation}
Let
\begin{equation}\label{eq:Bseries}
B(s,u):=\sum_{n\ge1}\frac{b_n(u)}{n^s}.
\end{equation}
In the initial half-plane of absolute convergence, Dirichlet multiplication gives
\begin{align}
B(s,u)
&=
\left(\sum_{k\ge1}\frac{e^{iu\gamma_k}}{k^s}\right)
\left(\sum_{d\ge1}\frac{\mu(d)}{d^{s+1}}\right)\notag\\
&=
\frac{Z_\gamma(s,u)}{\zeta(s+1)}.
\end{align}
Equivalently,
\begin{equation}\label{eq:threefactor}
\boxed{
Z_\gamma(s,u)=B(s,u)\,\zeta(s+1).
}
\end{equation}

The M\"obius--centroid construction supplies the holomorphic continuation of $B(s,u)$ across the left translated half-strip
\begin{equation}\label{eq:Bdomain}
\boxed{
-1<\Ree s< -\frac12,
\qquad \Im u>0.
}
\end{equation}
Thus, for every fixed $u$ with $\Im u>0$, $B(\,\cdot\,,u)$ is holomorphic at every point of \eqref{eq:Bdomain}.  Together with Proposition~\ref{prop:entire}, the identity \eqref{eq:threefactor} extends there by analytic continuation.

\section{Zero-free left half-strip}

\begin{theorem}[Forced zero-free translated strip]\label{thm:zerofree}
For every
\[
s_0\in\mathbb C,
\qquad
-1<\Ree s_0< -\frac12,
\]
one has
\[
\boxed{\zeta(s_0+1)\neq0.}
\]
\end{theorem}

\begin{proof}
Fix such an $s_0$.  By Lemma~\ref{lem:noncancel}, choose $y>Y(s_0)$ so that
\[
Z_\gamma(s_0,iy)\neq0.
\]
By the holomorphic construction of $B$ in \eqref{eq:Bdomain}, the value $B(s_0,iy)$ is finite.  Evaluating \eqref{eq:threefactor} at $(s_0,iy)$ gives
\[
Z_\gamma(s_0,iy)
=
B(s_0,iy)\,\zeta(s_0+1).
\]
If $\zeta(s_0+1)=0$, the right-hand side would vanish, contradicting
$Z_\gamma(s_0,iy)\neq0$.  Hence $\zeta(s_0+1)\neq0$.
\end{proof}

\begin{remark}[Multiplicity]
The argument does not assume simple zeros.  If $\zeta(s+1)$ had a zero of multiplicity $r\ge1$ at $s=s_0$, then the holomorphy of $B(s,iy)$ would force $Z_\gamma(s,iy)$ to vanish to order at least $r$ at $s_0$ through \eqref{eq:threefactor}.  Lemma~\ref{lem:noncancel} gives a value of $y$ for which already the zeroth-order value $Z_\gamma(s_0,iy)$ is nonzero, so every multiplicity $r\ge1$ is excluded at once.
\end{remark}

\section{Closure of the Riemann hypothesis}

\begin{theorem}[Riemann hypothesis]\label{thm:RH}
Every nontrivial zero $\rho$ of the Riemann zeta function satisfies
\[
\Ree\rho=\frac12.
\]
Hence $\RH$ holds.
\end{theorem}

\begin{proof}
Theorem~\ref{thm:zerofree} states that
\[
\zeta(s+1)\neq0
\qquad
\left(-1<\Ree s< -\frac12\right).
\]
Writing
\[
\rho=s+1,
\]
this is exactly
\[
\zeta(\rho)\neq0
\qquad
\left(0<\Ree\rho<\frac12\right).
\]
Thus there are no nontrivial zeros strictly to the left of the critical line.

The functional equation gives the symmetry
\[
\rho\longmapsto1-\rho
\]
of the nontrivial zeros, with multiplicity preserved.  Consequently, a zero with
\[
\Ree\rho>\frac12
\]
would produce a reflected zero $1-\rho$ with
\[
0<\Ree(1-\rho)<\frac12,
\]
contradicting the zero-free region just proved.  Therefore no nontrivial zero lies strictly on either side of the critical line, and every nontrivial zero satisfies
\[
\boxed{\Ree\rho=\frac12.}
\]
\end{proof}

\section{Logical core}

The proof reduces to the following three already established structural facts:
\[
\boxed{
\begin{array}{c}
Z_\gamma(s,u)\ \text{is holomorphic in }s\ \text{for }\Im u>0,\\[1mm]
B(s,u)\ \text{is holomorphic by construction for }-1<\Ree s<-\tfrac12,\\[1mm]
Z_\gamma(s,u)=B(s,u)\zeta(s+1).
\end{array}
}
\]
For each point $s_0$ of that strip, the large-$y$ non-cancellation lemma supplies $u=iy$ with
\[
Z_\gamma(s_0,iy)\neq0.
\]
Therefore the third factor in \eqref{eq:threefactor} cannot vanish there:
\[
\zeta(s_0+1)\neq0.
\]
This gives the zero-free half of the critical strip, and functional-equation symmetry supplies the other half.

\end{document}
